% Chapter 2 - Theoretical Background

\section{Theoretical Background}\label{sec:theory}

This chapter presents the scientific and technical foundations upon which the system rests.
The topics are introduced in pipeline order: from the anatomy of the human hand and its
skeletal parameterisation, through the models and algorithms that transform raw webcam
output into normalised joint-angle data, to the communication and rendering layers
that close the loop in real time.
The final section addresses the \acrshort{ai}-assisted surface reconstruction step that
converts user-painted point clouds into meshes.

\subsection{Anatomy of the Human Hand and Skeletal Modelling}\label{subsec:anatomy}

The human hand contains 27 bones grouped into three regions.
The eight \emph{carpal} bones form the wrist foundation.
The five \emph{metacarpal} bones span the palm, each articulating proximally with
a carpal and distally with a finger.
The 14 \emph{phalanges} are distributed as two in the thumb (proximal and distal)
and three in each of the remaining four fingers (proximal, middle, distal).
Motion arises at the joints linking these bones: the carpometacarpal, metacarpophalangeal,
and interphalangeal joints. Each contributes one or two degrees of freedom (\acrshort{dof}),
for more than twenty across the hand as a whole.

\begin{figure}[H]
	\centering
	\captionsetup{justification=centering}
	\includegraphics[width=0.5\linewidth, frame]{img/Scheme_human_hand_bones-en.png}
	\caption{Bones of the human hand}
	\label{fig:hand}
\end{figure}

For computational purposes, the hand is represented as a \emph{kinematic chain}:
a rooted tree of rigid segments (bones) connected at rotational joints.
Each joint $k$ has a designated parent joint $\mathrm{pa}(k)$, and its global
position in world space is obtained by composing all ancestor transforms along
the chain from the root (wrist) outward to the fingertips.
This parent--child hierarchy is the data structure underlying both MediaPipe's
21-landmark skeleton and the MANO parametric model's 16-joint armature.

Two complementary numbering conventions are relevant to this thesis.
Google MediaPipe assigns indices 0--20 to landmarks as follows: index 0 is the wrist,
indices 1--4 are the thumb (CMC, MCP, IP, TIP), indices 5--8 the index finger
(MCP, PIP, DIP, TIP), and so forth through to the little finger (indices 17--20).

\begin{figure}[H]
	\centering
	\captionsetup{justification=centering}
	\includegraphics[width=0.9\linewidth, frame]{img/MediaPipe-hand-model-Joint-annotation.png}
	\caption{MediaPipe hand model joint annotation. Adapted from Veluri et al. \cite{Veluri2022}.}
	\label{fig:mediapipe-joints}
\end{figure}

The MANO model, described in \autoref{subsec:mano}, uses a different 16-joint
convention in which each finger contributes three internal joints
(MCP, PIP, DIP) plus a fingertip keypoint extended from a specific mesh vertex.
Bridging these two representations is the central kinematic problem
addressed by the Adaptive Inverse Kinematics solver in \autoref{subsec:aik}.

\begin{figure}[H]
	\centering
	\captionsetup{justification=centering}
	\includegraphics[width=0.7\linewidth, frame]{img/MANO-hand-model-Joint-annotation.png}
	\caption{MANO hand model and index point labeling. Adapted from Ning et al. \cite{Ning2025}.}
	\label{fig:mano-joints}
\end{figure}
